\appendix

\chapter{锚点抽取提示词}
\label{app:prompt}

本节及附录 F 中的仓库路径均相对于 GitHub 仓库\url{\finplanrepo}。本节给出锚点抽取提示词的要点，完整提示词位于\repolink{experiments/anchor_cli/prompts/anchor_general.txt}{\texttt{experiments/anchor\_cli/prompts/anchor\_general.txt}}。

提示词的第一部分给出断言契约，包括公司、代码、目标期间、指标、截止时点、参考口径、结果标签及其定义。标签定义必须互斥且穷尽，例如“低于140\%”“140\%至170\%”“高于170\%”。提示词明确说明，指引区间定义三个标签，但区间本身不是“结果落在区间内”的证据。

第二部分给出时点证据。每条证据带来源编号$W\#$、标题、发布日期、链接和摘录。对于申报文件，摘录带段落编号$P\#$；对于结构化证据，编号为$E\#$。提示词要求锚点只引用这些编号，不引用不存在的来源。

第三部分给出任务和规则。规则包括：每条锚点必须引用一个来源编号；引文必须逐字复制；发布日期不得晚于截止时点；可靠性是隐含标签正确的校准概率；硬锚点必须是显式约束；监管截止日是上界，不是晚披露的证据；引文中提到的日期不得晚于来源发布日期；证据不足时在缺失字段中说明；标签分数必须使用契约标签且和为1。

第四部分给出输出格式。输出是一个JSON对象，包含锚点列表、标签分数、日期估计（仅日期断言）、置信度和缺失说明。锚点字段包括类型、关系、取值、作用域、来源编号、引文、日期、可靠性、硬标记、隐含标签和理由。日期估计字段包括点估计、区间下界、区间上界和依据。

提示词的设计目标是可检查性，而不是最大化抽取数量。系统宁可返回较少的锚点，也不接受没有来源或没有引文的锚点。解析失败的响应被记录为错误，不进入信念更新；系统不使用金标结果作为回退。

\chapter{数据构造与标签核验}
\label{app:data}

收入指引断言来自2010年至2025年美国SEC 8-K文件。每条断言包含以下字段：公司代码、公司名称、行业、目标年份、目标季度、指引下限、指引中点、指引上限、指引宽度、当前季度值、当前季度可靠性、截止时点、证据项、已实现结果和结果标签。

指引区间从8-K的业绩发布文本中抽取，单位为百万美元。当前季度值来自同一份8-K或最近的10-Q。已实现结果来自SEC XBRL；XBRL没有独立季度数据时，使用披露文件回退。结果标签按实际值相对指引区间的位置确定：低于下限为“低于”，落在区间内为“区间内”，高于上限为“高于”。

标签核验分为两级。第一级是XBRL核验：实际值来自XBRL的季度事实，且单位、期间和公司代码一致。第二级是XBRL覆盖：实际值来自XBRL覆盖的披露文件，但季度划分需要回退。227条测试断言满足其中一级；46条只有纯文本标签，被排除在主结果之外，用于敏感性检验。

A股回测覆盖40家非金融上市公司和2023年至2025年三个财年，共470条断言。财务数值来自akshare的公开财务报表，公告日期和链接来自巨潮资讯网。对每个财年$Y$，截止时点为$Y+1$年1月31日；营业收入增速、归母净利润增速和综合毛利率的阈值由上年实际值加减固定带宽得到，披露日期按3月31日和4月15日划分为三个窗口。只有上年值、三季报值、当年实际值和公告记录都存在的样本才进入基准。

数据的已知限制包括：标签来自自动核验而非人工金标准；类别不平衡；主基准只覆盖美国公司的季度营业收入指引；指引区间由公司选择，区间宽度与信息环境相关；部分证据项的日期来自文件提交日，可能与实际披露时点相差一天。这些限制在第\ref{ch:conclusion}章讨论。

\chapter{审计报告示例}
\label{app:audit}

本节给出中际旭创案例的审计报告结构。审计报告的生成方式与字段说明见\repolink{docs/anchor_cli.md}{\texttt{docs/anchor\_cli.md}}；运行后报告保存在本地结果目录的\texttt{audit.md}和\texttt{audit.json}中。

审计报告的第一部分是运行元数据：运行编号、生成时间、截止时点、问题、公司、代码、模型和配置。第二部分是结论摘要：每条断言的动作、点答案、裁决集、期望损失和披露日期估计。第三部分是断言契约：指标、类型、参考口径、结果标签及其定义。

第四部分是时点证据表：编号、日期、来源、标题、链接、检索词。官方公告记录排在前面，网页结果排在后面。第五部分是锚点与资格判定表：类型、取值或引文、来源、日期、引文核验、可靠性、资格、隐含标签和理由。被拒绝的锚点保留在表中，理由逐条列出。

第六部分是信念更新表：每个标签的先验、仅锚点后验和融合后验。第七部分是全动作损失感知读出：动作、点答案、裁决集、可信集、期望损失和各动作损失排序。第八部分是披露日期估计：点估计、区间、最高信念窗口、模型估计与窗口是否一致、依据。

审计报告的最后是审计说明：每条被采纳的锚点都有来源、引文、日期、资格和可靠性；每条被拒绝的锚点都有理由；晚于截止时点的证据被列入剔除列表；引文核验采用规范化后的逐字匹配；未通过核验的锚点被降权且不能成为硬锚点。报告明确说明，它是时点判断，不是对已实现结果的复核。

\chapter{关键比较的配对区间}
\label{app:paired}

表\ref{tab:paired}汇总正文关键比较的配对损失差与95\%置信区间。差值定义为前者减后者，负值表示前者损失更低。区间使用10,000次断言级配对自助法。

\begin{table}[htbp]
\centering
\zihao{5}
\setlength{\tabcolsep}{4pt}
\caption{关键比较的配对损失差与95\%置信区间（227条核验断言）。}
\label{tab:paired}
\begin{tabular}{llrr}
\toprule
比较 & 基座 & $\Delta$损失 & 95\%置信区间 \\
\midrule
全动作读出 $-$ 受限门 & DeepSeek & $-0.749$ & $[-1.075,\,-0.405]$ \\
 & GLM & $-0.581$ & $[-0.899,\,-0.264]$ \\
 & Qwen3-8B & $-0.696$ & $[-0.987,\,-0.405]$ \\
\midrule
融合 $-$ 随机森林 & DeepSeek & $-0.079$ & $[-0.300,\,0.141]$ \\
 & GLM & $-0.044$ & $[-0.203,\,0.115]$ \\
 & Qwen3-8B & $-0.062$ & $[-0.282,\,0.159]$ \\
\midrule
融合 $-$ 随机森林 & DeepSeek，指导区间内 & $-0.366$ & $[-0.732,\,-0.049]$ \\
 & GLM，指导区间内 & $-0.244$ & $[-0.537,\,-0.024]$ \\
 & Qwen3-8B，指导区间内 & $-0.415$ & $[-0.732,\,-0.170]$ \\
\bottomrule
\end{tabular}
\end{table}

\chapter{案例研究：中际旭创2023与2024年年报}
\label{app:case}

本案例使用已经发生的历史结果。对2023财年，截止时点为2024年1月31日，此时2023年三季报已经披露、2023年年报尚未披露；证据包括2021年和2022年年报、2023年三季报和2022年三季报。2023年年报于2024年4月22日披露，实际值用于验证预测。对2024财年，截止时点为2025年1月31日，此时2024年三季报已经披露、2024年年报尚未披露；2024年三季报和2023年年报进入证据，2024年年报于2025年4月21日披露。两个截止时点使用同一套契约、阈值规则、锚点抽取提示词和全动作读出。

表\ref{tab:case2023}给出2023年的运行结果。四条断言中，营业收入、综合毛利率和披露日期的点答案正确，归母净利润的点答案错误；三个点答案和三个裁决集覆盖了真实结果，归母净利润的真实值高于54.57\%，落在裁决集之外。归母净利润的数值桥估计为41.18\%，实际为77.58\%，误差主要来自2023年第四季度光模块需求的超预期增长，历史区间和前三季度年化都没有覆盖这一变化。这个错误保留在案例中，用来说明数值桥在需求突变时的失效方式。四条断言的后验按标签顺序分别为：营业收入 .607/.327/.066，归母净利润 .603/.331/.066，综合毛利率 .045/.420/.536，披露日期 .188/.223/.588。

\begin{table}[htbp]
\centering
\zihao{-5}
\setlength{\tabcolsep}{2pt}
\caption{中际旭创2023年年报案例（截止2024-01-31）。}
\label{tab:case2023}
\begin{tabular}{>{\raggedright\arraybackslash}p{3.0cm}>{\raggedright\arraybackslash}p{1.1cm}>{\raggedright\arraybackslash}p{2.6cm}>{\raggedright\arraybackslash}p{4.6cm}>{\raggedright\arraybackslash}p{2.0cm}}
\toprule
断言 & 动作 & 点答案 & 裁决集 & 实际 \\
\midrule
营业收入同比增速 & 裁决 & 低于15.29\% & 低于15.29\%, 15.29\%至35.29\% & 11.16\% \\
归母净利润同比增速 & 裁决 & 低于24.57\% & 低于24.57\%, 24.57\%至54.57\% & 77.58\% \\
综合毛利率 & 裁决 & 高于30.31pp & 28.31pp至30.31pp, 高于30.31pp & 32.99\% \\
披露日期 & 裁决 & 晚于2024-04-15 & 2024-04-01至04-15, 晚于04-15 & 2024-04-22 \\
\bottomrule
\end{tabular}
\end{table}

表\ref{tab:case2024}给出2024年的运行结果。2023年年报披露后，2023年实际值成为2024年断言的上年实际值，历史窗口从2021至2023年移动到2022至2024年，2024年三季报进入最新证据。四条断言的集合全部覆盖真实结果，点答案也全部正确。营业收入和归母净利润的数值桥估计分别为115.38\%和130.22\%，实际为122.64\%和137.93\%，误差比2023年明显缩小；综合毛利率的点答案为31.99pp至33.99pp，实际为33.80\%。

\begin{table}[htbp]
\centering
\zihao{-5}
\setlength{\tabcolsep}{2pt}
\caption{中际旭创2024年年报案例（截止2025-01-31）。}
\label{tab:case2024}
\begin{tabular}{>{\raggedright\arraybackslash}p{3.0cm}>{\raggedright\arraybackslash}p{1.1cm}>{\raggedright\arraybackslash}p{2.6cm}>{\raggedright\arraybackslash}p{4.6cm}>{\raggedright\arraybackslash}p{2.0cm}}
\toprule
断言 & 动作 & 点答案 & 裁决集 & 实际 \\
\midrule
营业收入同比增速 & 裁决 & 高于21.16\% & 低于1.162\%, 高于21.16\% & 122.64\% \\
归母净利润同比增速 & 裁决 & 高于92.58\% & 低于62.58\%, 高于92.58\% & 137.93\% \\
综合毛利率 & 点答 & 31.99pp至33.99pp & 31.99pp至33.99pp & 33.80\% \\
披露日期 & 裁决 & 晚于2025-04-15 & 2025-04-01至04-15, 晚于04-15 & 2025-04-21 \\
\bottomrule
\end{tabular}
\end{table}

{\sloppy 四条断言的后验按标签顺序分别为：营业收入 .188/.114/.698，归母净利润 .189/.117/.694，综合毛利率 .054/.818/.129，披露日期 .185/.220/.595。\par}

两个年度的对照展示了框架的演化过程。2023年结果揭示后，已实现值进入下一轮的上年实际值和历史窗口，锚点类型可靠性按揭示日更新，2024年断言的后验和动作随之变化。演化不是重写测试标签，而是把已经公开的结果放进下一时点的证据状态。2023年归母净利润的错误没有被掩盖，2024年同一指标的集合覆盖和点答案都得到修正；这比只展示成功案例更能说明框架在真实数据上的行为和边界。

案例的证据来自巨潮资讯网的年度报告和季度报告公告，财务数值来自公开财务报表。审计文件记录每条证据的来源、日期、引文、锚点资格和动作损失。案例的截止时点、阈值和标签定义在运行前固定，实际结果只在预测完成后用于验证。

\chapter{本地命令行与Claude Code接口}
\label{app:cli}

{\sloppy 本文实现了一个本地命令行工具和一个MCP服务。命令行工具为\repolink{scripts/anchor_cli.py}{\texttt{scripts/anchor\_cli.py}}，MCP服务为\repolink{scripts/anchor_mcp.py}{\texttt{scripts/anchor\_mcp.py}}，仓库根目录的\repolink{.mcp.json}{\texttt{.mcp.json}}注册了\texttt{finplan-anchor}服务。终端运行和Claude Code运行调用同一个引擎，生成同样的审计文件。\par}

命令行提供五个命令。\texttt{case}运行\texttt{configs/anchor\_cases}下的预置案例；\texttt{ask}把自然语言问题解析成断言契约并运行；\texttt{search}只检索证据，不做推断；\texttt{audit}打印某次运行的审计报告；\texttt{serve}启动MCP服务。典型用法为：

\begin{lstlisting}
git clone https://github.com/Topps-2025/FinPlanRAG.git
cd FinPlanRAG
python scripts/anchor_cli.py ask "中际旭创2023年年报增速区间" \
    --company 中际旭创 --ticker 300308.SZ \
    --period 2023年年度报告 --cutoff 2024-01-31
python scripts/anchor_cli.py ask "中际旭创2024年年报增速区间" \
    --company 中际旭创 --ticker 300308.SZ \
    --period 2024年年度报告 --cutoff 2025-01-31
python scripts/anchor_cli.py audit <run_id>
\end{lstlisting}

{\sloppy
MCP服务提供五个工具。\texttt{anchor\_ask}对自然语言问题做完整裁决；\texttt{anchor\_case}运行具名案例；\texttt{anchor\_search}检索带日期的证据；\texttt{anchor\_audit}读取审计报告；\texttt{anchor\_contract}只解析断言契约，便于用户先检查阈值和结果空间。Claude Code通过项目根目录的\texttt{.mcp.json}连接服务，用户可以直接说“用anchor\_case运行zhongji\_2026，并解释每个断言的证据链”。\par
}

接口的安全边界如下。API密钥只从环境变量读取，不写入代码、日志或审计文件。下载的网页和PDF作为不可信文本处理，不执行其中的脚本。LLM响应按请求负载哈希缓存，缓存目录不进入公开审计文件。审计文件保留证据链接和原文引文，便于复核，但不包含API密钥或用户隐私。

接口的局限如下。网页检索结果会随时间和搜索引擎状态变化，不同运行日的证据集合可能不同。审计文件记录当次实际使用的证据，不能把某一次运行的后验当作固定输出。案例文件中的阈值由用户或案例作者给定，阈值不合理时，框架只能在给定契约内做一致的推断。接口是研究框架的演示，不是交易系统，也不构成投资建议。
